\begin{table}[H]
\raggedright
\caption{Baseline characteristics of the analytic sample.}
\label{tab:sample}
\footnotesize
\begingroup
\setlength{\tabcolsep}{4pt}
\renewcommand{\arraystretch}{1.08}
\noindent
\begin{tabular*}{\textwidth}{@{\extracolsep{\fill}}lccc@{}}
\toprule
Variable & Mean (SD) & Range & n \\
\midrule
Age (years) & 46.4 (9.8) & 21 to 65 & 67 \\
Sex: women / men (\%) & 80.6 / 19.4 &  & 67 \\
Pain duration (months) & 82.8 (121.4) & 2 to 520 & 62 \\
MPI pain severity & 4.3 (0.9) & 2 to 6 & 67 \\
MPI interference & 4.2 (1.0) & 2 to 6 & 67 \\
Pain Disability Index & 42.7 (11.8) & 11 to 68 & 67 \\
PCS catastrophizing (total) & 22.3 (10.2) & 1 to 46 & 66 \\
PVAQ pain vigilance (total) & 43.7 (12.8) & 5 to 66 & 65 \\
HADS anxiety & 7.7 (3.9) & 1 to 20 & 66 \\
HADS depression & 8.7 (3.7) & 1 to 16 & 66 \\
NEO-FFI neuroticism & 30.8 (9.4) & 12 to 54 & 64 \\
\bottomrule
\end{tabular*}
\par\vspace{3pt}\noindent\parbox{\textwidth}{\footnotesize\textit{Note.} MPI = West Haven-Yale Multidimensional Pain Inventory; PDI = Pain Disability Index; PCS = Pain Catastrophizing Scale; PVAQ = Pain Vigilance and Awareness Questionnaire; HADS = Hospital Anxiety and Depression Scale; NEO-FFI = NEO Five-Factor Inventory.}
\endgroup
\end{table}
